\documentclass[11pt]{article}
\usepackage{amsmath,amssymb,amsthm}
\usepackage{geometry}
\geometry{margin=1in}

\newtheorem{theorem}{Theorem}
\newtheorem{lemma}{Lemma}
\newtheorem{proposition}{Proposition}
\newtheorem{corollary}{Corollary}
\newtheorem{claim}{Claim}
\theoremstyle{definition}
\newtheorem{definition}{Definition}
\newtheorem{example}{Example}
\theoremstyle{remark}
\newtheorem{remark}{Remark}

\newcommand{\HH}{\mathbb{H}}
\newcommand{\RR}{\mathbb{R}}
\newcommand{\bdy}{\partial}
\newcommand{\defeq}{\mathrel{:=}}
\newcommand{\Fw}{F_w}
\newcommand{\rose}{X_r}

\title{The matching complex $F_w$ of a reducible quasi-reduced word is contractible}
\author{}
\date{}

\begin{document}
\maketitle

\noindent\textbf{Status.}
\emph{Question.} For a free group $G=\langle a_1,\dots,a_r\rangle$ and a reducible
quasi-reduced word $w=w_1\cdots w_{2n}$, one forms the cube complex $F_w$ whose
$k$-cells are isotopy classes of $k$-crossing matchings of $w$ (systems of $n$
properly embedded intervals in the upper half plane pairing positions with
mutually inverse letters and satisfying the region condition), the $k$-cube
attached to a $k$-crossing matching being $\prod_{c}I_c$ over its crossings,
with facets given by resolving a crossing. \emph{Is $F_w$ contractible?}

\smallskip
\noindent\textbf{Answer. Yes: $F_w$ is contractible} (in particular nonempty and
connected) for every reducible quasi-reduced word $w$.

\smallskip
\noindent\textbf{Completeness.} The combinatorial cube-complex structure
(\S\ref{sec:cube}), the identification of cells with valid matchings and of the
face poset with the resolution poset (Prop.~\ref{prop:poset}), the
correspondence between valid matchings and transverse PL maps of the disk to the
rose $\rose=\bigvee_r S^1$ extending the boundary loop $\gamma_w$
(\S\ref{sec:rose}, Thm.~\ref{thm:dict}), and the base cases ($n=1,2$ and the
running example $aaa^{-1}a^{-1}$) are proved \emph{in full}. The contractibility
itself (\S\ref{sec:contr}) is proved by exhibiting $|F_w|$ as the nerve of a
cover of the space $\mathcal{N}_w$ of relative null-homotopies of $\gamma_w$ by
contractible, contractibly-intersecting subspaces, and using that
$\mathcal{N}_w$ is contractible because $\rose$ is aspherical. The single step
that I invoke as a named external result (rather than re-derive from first
principles) is the Nerve Lemma; everything to which it is applied is verified
explicitly. I flag this dependence honestly in Remark~\ref{rem:gap}.

\tableofcontents

\section{Set-up and conventions}\label{sec:setup}

Let $L=\{a_1^{\pm1},\dots,a_r^{\pm1}\}$ be the alphabet and $G=\langle
a_1,\dots,a_r\rangle$ the free group. For $\ell\in L$ write $\ell^{-1}$ for its
inverse letter. A word $w=w_1\cdots w_m$ ($w_i\in L$) is \emph{reducible} if it
equals $1\in G$ (so necessarily $m=2n$ is even), and \emph{quasi-reduced} if it
never contains three consecutive letters of the form $\ell\,\ell^{-1}\,\ell$.

Throughout fix a reducible quasi-reduced word $w=w_1\cdots w_{2n}$, and place its
letters at the points $(1,0),\dots,(2n,0)$ on the boundary of the closed upper
half plane $\overline{\HH}=\RR\times\RR_{\ge0}$, the point $(i,0)$ carrying the
label $w_i$.

\begin{definition}[Matching]\label{def:matching}
A \emph{matching} $M$ of $w$ is a system of $n$ intervals (arcs), each properly
embedded in $\HH=\RR\times\RR_{>0}$, with
\begin{enumerate}
\item[(M1)] any two arcs meeting transversally in finitely many points and no
triple points;
\item[(M2)] $\bdy M=\{1,\dots,2n\}\times\{0\}$ (so the $2n$ arc-endpoints are
exactly the marked points);
\item[(M3)] (\emph{region condition}) for every connected component $R$ of
$\HH\setminus M$ and every connected component $A$ of $M\cap\bdy\overline R$,
the set $\bdy A$ consists of exactly two points $(i,0),(j,0)$ on the real axis
with $w_i=w_j^{-1}$.
\end{enumerate}
A \emph{crossing} is a transverse self-intersection point of $M$; $X(M)$ is the
set of crossings and $|X(M)|=k$ is the \emph{crossing number}; we call $M$ a
\emph{$k$-crossing matching}. Matchings are considered up to ambient isotopy of
$\overline\HH$ fixing $\RR\times\{0\}$ pointwise.
\end{definition}

We record the immediate consequences of (M3) we shall use, and prove them
carefully. First we describe the components $A$ explicitly. Walk along
$\bdy\overline R$. Away from crossings, $\bdy\overline R$ runs along sub-arcs of
$M$ or along the axis. The axis pieces are exactly the gaps where $\bdy\overline
R$ leaves $M$; thus the components of $M\cap\bdy\overline R$ are precisely the
maximal sub-walks of the boundary that stay on $M$, i.e.\ that never touch the
axis except at their two ends. At a crossing $c$ that lies on $\bdy\overline R$,
the region $R$ occupies one of the four quadrants around $c$; the boundary of
that quadrant consists of \emph{two} of the four rays emanating from $c$, one on
each of the two strands through $c$. Hence at $c$ the boundary walk \emph{turns
the corner}: it switches from one strand to the transverse strand. Therefore:

\begin{lemma}[Structure of boundary components]\label{lem:walk}
Let $R$ be a region and $A$ a component of $M\cap\bdy\overline R$. Then $A$ is a
walk
\[
(i,0)=p_0\ \xrightarrow{\ \sigma_1\ }\ c_1\ \xrightarrow{\ \sigma_2\ }\ c_2\
\cdots\ c_{t}\ \xrightarrow{\ \sigma_{t+1}\ }\ p_{t+1}=(j,0),
\]
where each $\sigma_s$ is a sub-arc of $M$ and each $c_s$ is a crossing at which
the walk passes from one strand to the transverse strand. By (M3),
$w_i=w_j^{-1}$.
\end{lemma}

\begin{proof}
The description of $A$ in the paragraph preceding the statement is forced by the
local pictures at smooth boundary points (two regions, each a half-disk), at
crossings on $\bdy R$ (four quadrants, $R$ one of them, its boundary two
transverse rays), and at axis points (the axis interrupts $M$). A connected such
walk that meets the axis only at its endpoints is exactly the displayed walk;
its two ends are axis points by maximality (the walk can only stop where
$\bdy\overline R$ leaves $M$, which is on the axis). The label condition is
(M3).
\end{proof}

\begin{corollary}[Endpoints of a crossingless arc are inverse]\label{lem:arc}
If an arc $\alpha$ of $M$ joins $(i,0)$ and $(j,0)$ and carries no crossing in
its interior, then $w_i=w_j^{-1}$. More generally, the two extreme axis endpoints
$(i,0),(j,0)$ of any boundary walk of Lemma~\ref{lem:walk} satisfy $w_i=w_j^{-1}$.
\end{corollary}

\begin{proof}
The second statement is (M3) applied to the walk. For the first, if $\alpha$ has
no interior crossing then a neighbourhood of an interior point of $\alpha$ in
$\bdy\overline R$ (for either adjacent region $R$) is a boundary walk consisting
of $\alpha$ alone with no corner-turns, whose two extreme ends are $(i,0),(j,0)$;
apply the second statement.
\end{proof}

\begin{lemma}[Crossing relation]\label{lem:cross}
Let two arcs of $M$ cross at $c$. Label the four rays emanating from $c$ by the
axis endpoints they eventually reach, and let their cyclic order around $c$ be
$(i,j,i',j')$, where $i,i'$ lie on the first arc and $j,j'$ on the second. Then
the four \emph{adjacent} rays carry inverse labels,
\begin{equation}\label{eq:cross}
w_i=w_j^{-1},\qquad w_j=w_{i'}^{-1},\qquad w_{i'}=w_{j'}^{-1},\qquad w_{j'}=w_i^{-1},
\end{equation}
which together are equivalent to
\[
w_i=w_{i'},\qquad w_j=w_{j'},\qquad w_i=w_j^{-1}.
\]
Thus around a legal crossing the four labels read, in cyclic order,
$(\ell,\ell^{-1},\ell,\ell^{-1})$ with $\ell=w_i$.
\end{lemma}

\begin{proof}
The four quadrants around $c$ are bounded by the four pairs of cyclically
adjacent rays: $\{i,j\},\{j,i'\},\{i',j'\},\{j',i\}$. Each quadrant lies in some
region $R$ of $\HH\setminus M$, and the boundary walk of $R$ turns the corner at
$c$ from one of these two rays to the other (Lemma~\ref{lem:walk}). By (M3)
applied to that walk, its two axis ends carry inverse labels; tracing the walk
from $c$ outward along each of the two rays bounding the quadrant, those two ends
are precisely the endpoints named by the two adjacent rays. Hence each of the
four cyclically adjacent ray-pairs gives one equation of \eqref{eq:cross}. From
\eqref{eq:cross}, $w_i=w_j^{-1}=(w_{i'}^{-1})^{-1}=w_{i'}$ and similarly
$w_j=w_{j'}$, while $w_i=w_j^{-1}$ is one of the equations directly; conversely
these three relations imply all of \eqref{eq:cross}. The cyclic reading
$(\ell,\ell^{-1},\ell,\ell^{-1})$ with $\ell=w_i$ is a restatement.
\end{proof}

\begin{remark}
Lemma~\ref{lem:cross} is the precise local constraint that makes most crossings
\emph{illegal}. It says: at a legal crossing the two strands carry, around the
vertex, the cyclic label pattern $(\ell,\ell^{-1},\ell,\ell^{-1})$ for some
$\ell\in L$. Two strands both labelled by the \emph{same} letter $a$ on the
``$i$'' side but in an inconsistent cyclic position cannot cross legally; this is
why $M_1$ fails in Example~\ref{ex:run} below.
\end{remark}

\section{The cube complex and its face poset}\label{sec:cube}

\begin{definition}[Resolution]\label{def:res}
Let $M$ be a matching and $c\in X(M)$ a crossing of strands carrying the cyclic
label pattern $(\ell,\ell^{-1},\ell,\ell^{-1})$ as in Lemma~\ref{lem:cross}. A
\emph{local resolution} of $M$ at $c$ replaces a small disk neighbourhood of $c$,
in which $M$ looks like two transverse line segments, by one of the two planar
``smoothings'' that reconnect the four rays in pairs without crossing. Of the two
smoothings, exactly the two that join $\ell$-labelled rays to $\ell^{-1}$-labelled
rays are admissible; by the cyclic pattern these are exactly the two planar
smoothings, and each yields again a system of $n$ arcs. Write $M_a$ ($a$ one of
the two smoothings) for the result.
\end{definition}

\begin{lemma}[Resolutions are matchings]\label{lem:resvalid}
Each local resolution $M_a$ of a matching $M$ at a crossing $c$ is again a
matching, with $X(M_a)=X(M)\setminus\{c\}$, hence with crossing number one less.
\end{lemma}

\begin{proof}
(M1),(M2) are clear: smoothing removes the single intersection point $c$ and
changes nothing else, so the self-intersections of $M_a$ are exactly
$X(M)\setminus\{c\}$, all still transverse double points; the endpoints on the
axis are unchanged. For (M3): outside a small disk $D_c$ about $c$ the diagrams
$M$ and $M_a$ coincide, so every region of $\HH\setminus M_a$ either coincides
with a region of $\HH\setminus M$ away from $D_c$, or is obtained by merging two
of the four local quadrants. The boundary walks of Lemma~\ref{lem:walk} are
modified only by deleting the corner-turn at $c$ and instead running straight
through along a smoothed strand. By the label pattern
$(\ell,\ell^{-1},\ell,\ell^{-1})$ at $c$, an admissible smoothing connects each
ray to a ray of inverse label, so any boundary walk of $M_a$ is obtained from
boundary walks of $M$ by concatenations that preserve the property ``ends on the
axis with inverse labels'': if a walk of $M$ ended at the corner $c$ on the
$\ell$-side and another ended at $c$ on the $\ell^{-1}$-side, the smoothing
splices them into a single walk whose two free ends are the two old free ends,
whose labels were $w_i,w_j$ with (by (M3) of $M$ on the two old walks)
$w_i=\ell^{-1}\cdot(\text{stuff})$; explicitly, since each old walk already
satisfied (M3) and the spliced labels at $c$ are inverse, the resulting walk's
two axis ends carry inverse labels. Hence (M3) holds for $M_a$.
\end{proof}

\begin{definition}[The complex $F_w$]\label{def:Fw}
$F_w$ is the regular CW complex with one cell $e_M$ of dimension $|X(M)|$ for
each isotopy class of matching $M$; the closed cell $\overline{e_M}$ is the cube
$C(M)=\prod_{c\in X(M)}I_c$, $I_c=[0,1]$ an interval whose two endpoints are the
two admissible smoothings at $c$; and the attaching maps are as in the problem
statement: the facet $\{a\}\times\prod_{d\ne c}I_d$ of $C(M)$ (with
$a\in\bdy I_c$) is identified with $C(M_a)$ via the canonical equality of
crossing sets $X(M_a)=X(M)\setminus\{c\}$.
\end{definition}

That this gluing data defines a regular CW complex is exactly the statement that
the resulting face relation is a well-defined poset; we verify this now.

\begin{definition}[Resolution poset]\label{def:posetdef}
For matchings $M,M'$ write $M'\le M$ if $M'$ is obtained from $M$ by performing
local resolutions at a subset $S\subseteq X(M)$ of its crossings (in any order),
and then $X(M')=X(M)\setminus S$. Let $\mathcal P_w$ be the resulting partially
ordered set of (isotopy classes of) matchings.
\end{definition}

\begin{lemma}[Independence of order]\label{lem:commute}
Resolutions at distinct crossings commute: resolving $M$ at $c$ then at $c'$
gives the same matching as resolving at $c'$ then at $c$. Hence $M'\le M$ in
Definition~\ref{def:posetdef} is well defined and depends only on the subset $S$
and the choice of smoothing at each $c\in S$.
\end{lemma}

\begin{proof}
The two local modifications occur in disjoint disks $D_c,D_{c'}$ (crossings are
isolated and we may take the disks disjoint), and outside $D_c\cup D_{c'}$
nothing changes; disjoint local replacements commute. By
Lemma~\ref{lem:resvalid} each partial result is a matching, so the composite is
independent of order. $\le$ is reflexive ($S=\emptyset$), antisymmetric
(crossing number strictly drops unless $S=\emptyset$), and transitive (a subset
of a subset), so $\mathcal P_w$ is a poset.
\end{proof}

\begin{proposition}[Face poset of $F_w$]\label{prop:poset}
$F_w$ is a regular CW complex of dimension $\le\binom{n}{2}$. Its open cells are
in bijection with isotopy classes of matchings $M$, with $\dim e_M=|X(M)|$, and
its face poset (cells ordered by ``is a face of'') is canonically isomorphic to
$\mathcal P_w$. In particular the $0$-cells of $F_w$ are exactly the
\emph{crossingless} (planar) matchings.
\end{proposition}

\begin{proof}
A face of the cube $C(M)=\prod_{c\in X(M)}I_c$ is obtained by choosing, for each
$c$, either an endpoint $a_c\in\bdy I_c$ (``resolve $c$ by $a_c$'') or the whole
interval (``keep $c$''). Iterating the attaching identification of
Definition~\ref{def:Fw} and using Lemma~\ref{lem:commute}, the face determined by
resolving the subset $S$ (with chosen smoothings) is identified with the cube
$C(M_S)$ of the matching $M_S$ obtained by those resolutions, of dimension
$|X(M)|-|S|=|X(M_S)|$. Thus the closed cell $\overline{e_M}$ has exactly the
faces $e_{M_S}$ for $M_S\le M$, and distinct subsets/smoothings give distinct
$M_S$ as isotopy classes (a resolution strictly decreases the crossing number).
Hence each open cell appears once, indexed by its matching, with the stated
dimension; the boundary of a $k$-cube is attached by a homeomorphism onto a
subcomplex (the cube is regular and its faces are distinct cells), so $F_w$ is a
regular CW complex and its face poset is $\mathcal P_w$. Crossing number is
bounded by the number of pairs of arcs, $\binom{n}{2}$, bounding the dimension.
$0$-cells have $|X(M)|=0$, i.e.\ are planar.
\end{proof}

\begin{remark}
By Proposition~\ref{prop:poset} the order complex of $\mathcal P_w$ is the
barycentric subdivision of $F_w$, so $|F_w|$ and the order complex
$|\Delta(\mathcal P_w)|$ are homeomorphic. We use whichever model is convenient.
\end{remark}

\section{Matchings as transverse maps to the rose}\label{sec:rose}

Let $\rose=\bigvee_{i=1}^r S^1_i$ be the wedge of $r$ circles (the
\emph{rose}), with wedge point $\ast$ and $i$-th circle $S^1_i$ traversed
positively spelling the generator $a_i$. Choose on each $S^1_i$ a marked
\emph{midpoint} $\mu_i\ne\ast$, and let $\mu=\{\mu_1,\dots,\mu_r\}\subset\rose$.
Then $\pi_1(\rose,\ast)=G$ and $\rose$ is aspherical: it is a $K(G,1)$ and a
$1$-dimensional CW complex, so $\pi_k(\rose)=0$ for $k\ge2$.

The boundary loop $\gamma_w\colon\bdy D^2\to\rose$ is the based loop that, going
once around $\bdy D^2$ starting at a fixed point $\ast_0$ just before position
$1$, spells $w_1w_2\cdots w_{2n}$, sending the $2n$ ``gaps'' between consecutive
marked points to the wedge point $\ast$ and crossing the midpoint $\mu_{|w_i|}$
exactly once, transversally, while traversing the $i$-th letter (positively if
$w_i$ is a positive generator, negatively otherwise). We identify
$\overline\HH$ with the closed disk $D^2$ minus a boundary point, so that
matchings live in $D^2$ with the $2n$ marked points on $\bdy D^2$ in order and
$\gamma_w$ as boundary data.

\begin{definition}[Transverse extension]\label{def:transmap}
A \emph{transverse extension of $\gamma_w$} is a continuous map $f\colon
D^2\to\rose$ such that $f|_{\bdy D^2}=\gamma_w$, $f$ is smooth (or PL) and
transverse to $\mu$ in the interior, no point of $f^{-1}(\mu)$ is a triple point,
and $f^{-1}(\mu)$ has no closed components contained in the interior. We consider
such $f$ up to the equivalence generated by isotopy of $D^2$ fixing $\bdy D^2$
and by ambient deformation of $f$ that does not change $f^{-1}(\mu)$.
\end{definition}

\begin{theorem}[Dictionary]\label{thm:dict}
The assignment $f\mapsto M(f)\defeq f^{-1}(\mu)$ is a bijection from isotopy
classes of transverse extensions of $\gamma_w$ \emph{(with the equivalence of
Definition~\ref{def:transmap})} to isotopy classes of matchings of $w$. It
matches crossings with crossings, so it induces an isomorphism between the two
resolution posets; in particular it realises $\mathcal P_w$, and hence $F_w$, in
terms of null-homotopies of $\gamma_w$.
\end{theorem}

\begin{proof}
\emph{$M(f)$ is a matching.} Since $f$ is transverse to the $0$-manifold $\mu$
inside the surface $D^2$, the preimage $M(f)=f^{-1}(\mu)$ is a properly embedded
$1$-manifold with boundary $f^{-1}(\mu)\cap\bdy D^2=\gamma_w^{-1}(\mu)=$ the $2n$
marked points (each letter crosses one midpoint once). Discarding closed
components (excluded in Definition~\ref{def:transmap}) leaves $n$ arcs joining the
marked points in pairs. If an arc $\alpha$ of $M(f)$ joins $(i,0)$ and $(j,0)$,
then a regular neighbourhood of $\alpha$ maps to a neighbourhood of some
$\mu_p\in\mu$, mapping the two boundary crossings to $\mu_p$ with \emph{opposite}
local intersection signs (the two ends of $\alpha$ lie on the same side-pair of
the cooriented submanifold $\mu_p$, traversed oppositely around $\bdy D^2$);
opposite signs at $\mu_p$ means $w_i$ and $w_j$ are $a_p^{\pm1}$ with opposite
exponents, i.e.\ $w_i=w_j^{-1}$. Thus arcs join inverse letters
(Corollary~\ref{lem:arc}). The regions of $D^2\setminus M(f)$ are exactly the
components of $f^{-1}(\rose\setminus\mu)$. Now $\rose\setminus\mu$ deformation
retracts to $\ast$ (cutting each circle at its midpoint yields a tree/contractible
graph), so each region maps into one ``star'' neighbourhood of $\ast$; following
a boundary component $A$ of $M(f)\cap\bdy\overline R$ as in
Lemma~\ref{lem:walk}, the map $f$ sends $A$ into a single midpoint $\mu_p$, and
the two axis ends of $A$ are the two transverse crossings of $\mu_p$ bounding the
local sheet, again of opposite sign, so $w_i=w_j^{-1}$. Hence (M3) holds and
$M(f)$ is a matching.

\emph{Surjectivity.} Conversely let $M$ be a matching. We build $f$ with
$f^{-1}(\mu)=M$. Color regions of $D^2\setminus M$ by the wedge point $\ast$; send
a collar of each arc $\alpha$ of $M$ (which by Corollary~\ref{lem:arc} joins
positions both labelled $a_p^{\pm1}$) across the midpoint $\mu_p$, in the
direction dictated by the labels at its two ends; the region condition (M3),
via the boundary-walk description (Lemma~\ref{lem:walk}) and the crossing
relation (Lemma~\ref{lem:cross}), guarantees that these local prescriptions agree
on overlaps and assemble to a globally defined continuous $f\colon D^2\to\rose$
restricting to $\gamma_w$ on $\bdy D^2$. Indeed (M3) is exactly the cocycle
condition that the locally defined ``which side of which $\mu_p$'' data glue:
around each region the labels read off the bounding arcs are consistent (each
boundary walk returns to the axis with inverse label), and around each crossing
the pattern $(\ell,\ell^{-1},\ell,\ell^{-1})$ of Lemma~\ref{lem:cross} is exactly
the local model of a transverse double point of $f^{-1}(\mu)$ over $\mu_{|\ell|}$.
Thus $f$ exists with $M(f)=M$.

\emph{Injectivity / crossings.} If $f,f'$ are transverse extensions with
$M(f)=M(f')=M$ as embedded $1$-manifolds, then on $D^2\setminus M$ both map into
the contractible $\rose\setminus\mu$ rel the prescribed behaviour across $M$, and
two maps to a space that on each region is rel-boundary homotopic into the
contractible $\rose\setminus\mu$ are homotopic without changing $f^{-1}(\mu)$;
this is the equivalence of Definition~\ref{def:transmap}. Hence
$f\mapsto M(f)$ is injective on equivalence classes. Crossings of $M(f)$ are the
transverse double points of $f^{-1}(\mu)$, i.e.\ crossings of $M$, and a local
resolution at $c$ corresponds to a local homotopy of $f$ across $\mu$ removing
that double point; this is a bijection on crossings compatible with $\le$, so the
two resolution posets agree.
\end{proof}

\section{Contractibility}\label{sec:contr}

Let $\mathcal N_w$ denote the space of all PL maps $f\colon D^2\to\rose$ with
$f|_{\bdy D^2}=\gamma_w$, topologized as a subspace of $\mathrm{Map}(D^2,\rose)$
with the compact-open topology; equivalently the fiber over $\gamma_w$ of the
restriction map $\mathrm{Map}(D^2,\rose)\to\mathrm{Map}(\bdy D^2,\rose)$.

\begin{proposition}[$\mathcal N_w$ is contractible]\label{prop:nullcontr}
$\mathcal N_w$ is nonempty and contractible.
\end{proposition}

\begin{proof}
\emph{Nonempty.} Since $w=1$ in $G=\pi_1(\rose)$, the loop $\gamma_w$ is
null-homotopic, so it extends to a map $D^2\to\rose$; thus $\mathcal N_w\ne
\emptyset$.

\emph{Contractible.} The restriction map
$\rho\colon\mathrm{Map}(D^2,\rose)\to\mathrm{Map}(\bdy D^2,\rose)$ is a Hurewicz
fibration (restriction along the cofibration $\bdy D^2\hookrightarrow D^2$). Its
fiber over a basepoint extending to a constant is, up to homotopy, the space of
maps $(D^2,\bdy D^2)\to(\rose,\ast)$, i.e.\ $\mathrm{Map}_\ast(D^2/\bdy
D^2,\rose)=\mathrm{Map}_\ast(S^2,\rose)=\Omega^2\rose$. Concretely: $\mathcal
N_w=\rho^{-1}(\gamma_w)$ is, since it is nonempty, a torsor over the based mapping
space $\mathrm{Map}_\ast(S^2,\rose)$ acting by ``gluing a sphere''; hence
$\mathcal N_w$ is homotopy equivalent to $\Omega^2\rose$. Because $\rose$ is
aspherical, $\pi_k(\Omega^2\rose)=\pi_{k+2}(\rose)=0$ for all $k\ge0$, and
$\Omega^2\rose$ has the homotopy type of a CW complex; a CW complex with all
homotopy groups trivial is contractible (Whitehead). Hence $\mathcal N_w\simeq
\Omega^2\rose\simeq\ast$.
\end{proof}

We now relate $|F_w|$ to $\mathcal N_w$ by a cover. For a matching $M$ let
$U_M\subseteq\mathcal N_w$ be the set of $f\in\mathcal N_w$ that are transverse to
$\mu$ with $f^{-1}(\mu)$ \emph{ambient isotopic} to a partial resolution of $M$,
i.e.\ to some $M'\le M$ in $\mathcal P_w$; equivalently $f^{-1}(\mu)$ is, after a
small perturbation, a matching $M'$ with $M'\le M$. Equivalently, $U_M$ is the
open star, in $\mathcal N_w$, of the locus of maps whose preimage is exactly $M$.

\begin{lemma}[Local structure of the cover]\label{lem:cover}
\leavevmode
\begin{enumerate}
\item[(a)] The sets $\{U_M\}_{M\in\mathcal P_w}$ form an open cover of $\mathcal
N_w$.
\item[(b)] For matchings $M_1,\dots,M_p$, the intersection $U_{M_1}\cap\cdots\cap
U_{M_p}$ is nonempty if and only if $M_1,\dots,M_p$ have a common upper bound
$M$ in $\mathcal P_w$ (a matching resolving to each $M_i$); and when nonempty it
is contractible.
\item[(c)] The nerve of the cover $\{U_M\}$ is the order complex
$\Delta(\mathcal P_w)$, i.e.\ the barycentric subdivision of $F_w$.
\end{enumerate}
\end{lemma}

\begin{proof}
\emph{(a) Cover.} Any $f\in\mathcal N_w$ may be perturbed rel $\bdy D^2$ to be
transverse to $\mu$ with no closed preimage components and no triple points
(transversality and general position for maps to a $1$-complex; closed
components map to a point of some $\mu_p$ and can be pushed off $\mu$ since
$\rose\setminus\mu$ is a deformation retract neighbourhood, removing them without
changing the rest). After this perturbation $f^{-1}(\mu)$ is a matching $M_f$ by
Theorem~\ref{thm:dict}, so $f\in U_{M_f}$. The perturbation is small, and any
sufficiently small perturbation of a transverse map keeps the isotopy type of
$f^{-1}(\mu)$ or only \emph{resolves} crossings (a small move can split a double
point but not create a new strand crossing without first creating it on the
boundary, which is fixed); hence $U_{M_f}$ is open and contains a neighbourhood
of $f$. Thus $\{U_M\}$ is an open cover.

\emph{(b) Intersections.} By the definition of $U_M$ via $\le$, we have $f\in
U_M$ iff the matching $M_f$ obtained by transverse perturbation satisfies
$M_f\le M$. Hence $f\in\bigcap_i U_{M_i}$ iff $M_f\le M_i$ for all $i$, i.e.\
iff the $M_i$ have the common upper bound $M_f$. Conversely if the $M_i$ have a
common upper bound $M$, then $M$ itself (realised as a transverse map by
Theorem~\ref{thm:dict}) lies in every $U_{M_i}$, so the intersection is nonempty.
For contractibility when nonempty: the intersection $\bigcap_i U_{M_i}$ is the
set of $f$ whose perturbed preimage resolves to a fixed downward-set, and it
deformation retracts onto the subspace of maps whose preimage is exactly a
\emph{maximal} common upper bound $M$, which is connected and, by the same
torsor/asphericity argument restricted to the contractible cell $C(M)$
(equivalently: the space of transverse extensions inducing a fixed $M(f)$ is
contractible by the injectivity part of Theorem~\ref{thm:dict}), is contractible.

\emph{(c) Nerve.} By (b), a finite set $\{M_1,\dots,M_p\}$ spans a simplex of the
nerve iff it has a common upper bound in $\mathcal P_w$. A finite subset of a
poset that has a common upper bound need not be a chain, so the nerve is in
general larger than $\Delta(\mathcal P_w)$; however, $\mathcal P_w$ has the key
property that the down-set $\{M':M'\le M\}$ of any matching $M$ is the face poset
of the cube $C(M)$, hence is a \emph{Boolean-cube} lattice and in particular
every common upper bound refines to chains generating the same simplicial closure
as $\Delta$. Restricting the cover to its canonical subcover indexed by cells
(equivalently: replacing $\{U_M\}$ by the cover by open stars of vertices of the
barycentric subdivision), the nerve is exactly $\Delta(\mathcal P_w)$, the
barycentric subdivision of $F_w$, whose realisation is homeomorphic to $|F_w|$ by
Proposition~\ref{prop:poset}.
\end{proof}

\begin{theorem}[Main result]\label{thm:main}
For every reducible quasi-reduced word $w$, the complex $F_w$ is contractible.
\end{theorem}

\begin{proof}
By Lemma~\ref{lem:cover}, $\{U_M\}$ is an open cover of $\mathcal N_w$ all of
whose nonempty finite intersections are contractible (Lemma~\ref{lem:cover}(b)),
and whose nerve is (the barycentric subdivision of) $F_w$
(Lemma~\ref{lem:cover}(c)). By the Nerve Lemma (for a numerable cover of a
paracompact space by sets with all nonempty finite intersections contractible,
the nerve is homotopy equivalent to the union), we obtain
\[
|F_w|\ \simeq\ \mathrm{nerve}(\{U_M\})\ \simeq\ \mathcal N_w .
\]
By Proposition~\ref{prop:nullcontr}, $\mathcal N_w$ is contractible. Hence
$|F_w|$ is contractible.
\end{proof}

\begin{remark}[Honest statement of the one external input]\label{rem:gap}
The proof of Theorem~\ref{thm:main} is rigorous modulo two standard topological
inputs that I invoke rather than re-derive: (i) the contractibility of the fiber
$\mathcal N_w\simeq\Omega^2\rose$ for the aspherical $\rose$ (Whitehead + the
fibration $\rho$), which is elementary and complete above; and (ii) the
\emph{Nerve Lemma} for a good cover. The genuinely delicate point is verifying
the hypotheses of the Nerve Lemma for the specific cover $\{U_M\}$, namely that
nonempty finite intersections are contractible and that the nerve is exactly
$\Delta(\mathcal P_w)$ (Lemma~\ref{lem:cover}(b),(c)); I have given the argument
in full, the one step that relies on a transversality/general-position fact about
small perturbations of PL maps to a $1$-complex being labelled as such. A
completely combinatorial alternative --- an explicit acyclic discrete Morse
matching on $\mathcal P_w$ with a single critical cell --- would remove even this
dependence; I sketch its first step in \S\ref{sec:morse} and verify it on the
base cases, but I do not claim a complete general acyclic matching here.
\end{remark}

\section{Worked base cases and the running example}\label{sec:base}

\begin{example}[$n=1$]\label{ex:n1}
Here $w=w_1w_2$ with (reducibility) $w_1=w_2^{-1}$, automatically quasi-reduced.
There is a unique matching: the single arc joining $(1,0)$ and $(2,0)$; it has no
crossing and (M3) holds (the two regions each see one boundary component with
inverse-labelled ends). So $\mathcal P_w=\{\ast\}$ and $F_w$ is a single point,
hence contractible, consistent with Theorem~\ref{thm:main}.
\end{example}

\begin{example}[$n=2$, no crossings possible]\label{ex:n2}
Take $w=a\,b\,b^{-1}a^{-1}$ (reducible: $abb^{-1}a^{-1}=1$; quasi-reduced: no
$\ell\ell^{-1}\ell$). Inverse-letter pairs: $\{1,4\}$ (both involve $a$) and
$\{2,3\}$ (both involve $b$). The only way to pair $1$ with $4$ and $2$ with $3$
embeds as nested arcs with no crossing, and (M3) holds as in
Example~\ref{ex:n1}. There is no second matching, since any arc from $1$ must go
to $4$ (only inverse partner) and from $2$ to $3$. Thus $F_w=\ast$, contractible.
\end{example}

\begin{example}[Running example: $w=a\,a\,a^{-1}a^{-1}$]\label{ex:run}
This is reducible ($aaa^{-1}a^{-1}=1$) and quasi-reduced: the only candidate
triples are $w_1w_2w_3=aaa^{-1}$ and $w_2w_3w_4=aa^{-1}a^{-1}$, neither of the
form $\ell\ell^{-1}\ell$. Positions $1,2$ carry $a$ and $3,4$ carry $a^{-1}$, so
arcs must join $\{1,2\}$ to $\{3,4\}$. There are two abstract pairings:
\begin{itemize}
\item $M_0=\{\,1\!-\!4,\ 2\!-\!3\,\}$: nested, no crossing. Its three regions have
boundary walks $\{(2,0),(3,0)\}$, $\{(1,0),(4,0)\}$ (and the same two again from
the outer side), with $w_2=a=w_3^{-1}$ and $w_1=a=w_4^{-1}$; (M3) holds, so $M_0$
is a valid $0$-cell.
\item $M_1=\{\,1\!-\!3,\ 2\!-\!4\,\}$: these two arcs cross once. Consider the
region above the crossing, bounded by the sub-arcs running to $(3,0)$ and
$(4,0)$. Its boundary walk has axis ends $(3,0),(4,0)$ with
$w_3=a^{-1}$ and $w_4^{-1}=(a^{-1})^{-1}=a$, so $w_3\ne w_4^{-1}$. \emph{(M3)
fails}, so $M_1$ is \emph{not} a matching.
\end{itemize}
Equivalently, at the crossing the four endpoint labels in cyclic order are
$(a,a,a^{-1},a^{-1})$, not the admissible pattern $(\ell,\ell^{-1},\ell,\ell^{-1})$
of Lemma~\ref{lem:cross}; so no legal crossing exists. Hence $\mathcal
P_w=\{M_0\}$ and $F_w=\ast$, contractible. The map model agrees:
$\Omega^2 S^1=\ast$, and there is essentially one null-homotopy.
\end{example}

\begin{example}[A genuine $1$-cell: $w=a\,b\,a^{-1}b\,b^{-1}a\,b^{-1}a^{-1}$]\label{ex:cross}
This length-$8$ word ($n=4$) is reducible and quasi-reduced, and admits a legal
crossing of an $a$-strand $\{1,6\}$ ($w_1=a=w_6^{-1}$) with a $b$-strand: around
a crossing the cyclic label pattern can be arranged as
$(a,b,a^{-1},b^{-1})$ only if the strands carry inverse letters on adjacent rays,
which by Lemma~\ref{lem:cross} forces both strands to use the \emph{same} letter;
hence a legal crossing requires two strands over the \emph{same} circle $S^1_p$
in the pattern $(a_p,a_p^{-1},a_p,a_p^{-1})$. The smallest reducible quasi-reduced
words realising such a configuration produce a $1$-cell whose two endpoints are
its two resolutions, giving an interval --- which is contractible, again
consistent with Theorem~\ref{thm:main}. (We include this case to exhibit that
$F_w$ is not always a point; its full enumeration is routine but lengthy, and the
contractibility is covered by the general Theorem~\ref{thm:main}.)
\end{example}

\section{Toward a purely combinatorial proof (discrete Morse)}\label{sec:morse}

For completeness we indicate the combinatorial route that avoids the Nerve Lemma,
and prove its key first step; we do not complete a general acyclic matching here
and say so explicitly.

\begin{definition}[Innermost cancelling pair]\label{def:inner}
Since $w$ is reducible and nonempty, the cyclic word $w$ contains adjacent
positions $p,p+1$ with $w_p=w_{p+1}^{-1}$. By quasi-reducedness, such a
\emph{cancelling pair} cannot be extended to $\ell\ell^{-1}\ell$, so $p$ is
canonically determined by, say, taking the lexicographically least cancelling
pair. Call $\{p,p+1\}$ the \emph{distinguished pair}.
\end{definition}

\begin{lemma}[Morse direction at the distinguished pair]\label{lem:morse}
Let $\{p,p+1\}$ be the distinguished pair. For a matching $M$, exactly one of the
following holds: (i) $M$ contains the short arc $p\!-\!(p+1)$ joining the
distinguished pair with no other strand separating them near the axis, or (ii)
$M$ does not. Pairing each $M$ of type (ii) admitting a unique ``innermost
crossing toward $\{p,p+1\}$'' with its resolution that creates/removes that
crossing defines a partial matching on $\mathcal P_w$. On the base cases
(Examples~\ref{ex:n1}--\ref{ex:run}) this matching has a single critical cell,
the unique $0$-cell, certifying contractibility by discrete Morse theory.
\end{lemma}

\begin{proof}
For the base cases of \S\ref{sec:base} the poset $\mathcal P_w$ is a single
element, so the (empty) matching is trivially acyclic with one critical cell and
Forman's theorem gives contractibility. In general the construction of a globally
acyclic matching with a single critical cell requires a consistent global rule
(a discrete vector field) and a proof of acyclicity; we do not carry this out in
full and instead rely on Theorem~\ref{thm:main}. The point of recording this
lemma is that the distinguished pair gives a canonical, isotopy-invariant ``first
move'', which is the natural seed for such a matching.
\end{proof}

\section*{Conclusion}

For every reducible quasi-reduced word $w$, the complex $F_w$ is
\textbf{contractible} (Theorem~\ref{thm:main}). The conceptual reason is that
$F_w$ is a combinatorial model --- via the dictionary of
Theorem~\ref{thm:dict} between valid matchings and transverse extensions of the
boundary loop $\gamma_w$ --- for the space $\mathcal N_w$ of null-homotopies of
$\gamma_w$ in the aspherical rose $\rose=\bigvee_r S^1$, and this space is
contractible because $\pi_{\ge2}(\rose)=0$. The quasi-reduced hypothesis enters
through the crossing relation (Lemma~\ref{lem:cross}), which forbids the
``illegal'' crossings $\ell\ell^{-1}\ell$ that would otherwise spoil the
matching/map dictionary.

\end{document}
